%%%%%%%%%%%%%%%%%%%%%%%%%%%%%%%%%%%%%%%%%%%%%%%%%%%%%%%%%%%%%%%%%%
% admissions-plan.tex -- PhD Proposed Plan of Study (application document)
% York/TMU Communication & Culture, Fall 2027.
%
% York's instruction (How to Apply page, re-checked 2026-10-13):
% approximately 500 words; "a clear and detailed description of the planned
% research area by including the objects of study, potential methods, and
% the compatibility of their scholarly interests with the program", plus
% preparation for doctoral study where earlier background "appears distant"
% from communication and culture; relevant accomplishments may be noted.
% This is NOT the formal dissertation proposal (max. 3,500 words) that is
% prepared after admission with a supervisory committee.
%
% Layout follows the York thesis template supplied by the author; no logo.
% Compiles with pdfLaTeX, XeLaTeX or LuaLaTeX.
%%%%%%%%%%%%%%%%%%%%%%%%%%%%%%%%%%%%%%%%%%%%%%%%%%%%%%%%%%%%%%%%%%
\documentclass[letterpaper,11pt]{article}
\usepackage{iftex}
\ifPDFTeX
  \usepackage[T1]{fontenc}
  \usepackage[utf8]{inputenc}
  \usepackage{mathptmx}
\else
  \usepackage{fontspec}
  \setmainfont{texgyretermes}[Extension=.otf, UprightFont=*-regular,
    BoldFont=*-bold, ItalicFont=*-italic, BoldItalicFont=*-bolditalic]
\fi
\usepackage[top=25mm, bottom=25mm, left=25mm, right=25mm]{geometry}
\usepackage{setspace}
\usepackage[hidelinks]{hyperref}
\setlength{\parindent}{0pt}
\setlength{\parskip}{6pt}
\onehalfspacing
\pagestyle{empty}

\begin{document}

\begin{center}
  {\bf\large Proposed Plan of Study}\\[2pt]
  Roman Jahandideh --- PhD in Communication \& Culture, Fall 2027\\
  \textit{Whose Contribution Counts? AI Mediation and Collective Authorship in Digital Improvisation}
\end{center}

In collaborative creative tools, a generative model reinterprets what
participants contribute: one person's phrase becomes a sound or a moving form
in a shared scene, and its outputs may become more visible than the
contributions that prompted them. My doctoral research asks how participants, facilitators
and a generative-AI interface shape whether a person's contribution is
recognised, carries creative authority, and is credited in a collective work.
I will examine whether technical attribution, the system's
record of who supplied what, and social recognition among collaborators come
apart. The question concerns communication and cultural participation: whose
interpretations gain authority, and how authorship is negotiated when a
machine takes part in creative exchange.

My objects of study are small-group digital improvisation workshops with
adult improvising artists and creative practitioners, and the shared
audiovisual environment I will build for them. Participants contribute short
phrases. An existing language model proposes transformations into light,
form and sound, and contributors or groups can accept, revise, reject or
undo them. Anyone can take part without using AI.

I plan to use the program's research-creation option. The environment and
the improvised works made with it form the creative component, and
developing its rules and workshop scores is part of the inquiry. The methods
are:
\begin{itemize}\setlength{\itemsep}{0pt}
  \item scoping interviews with improvisers and facilitators;
  \item co-design with practitioners;
  \item repeated workshops (provisionally three small groups, four sessions
        each) comparing contributor-led, group-led and human-mediated
        arrangements;
  \item reflection interviews in which participants review video of their
        own sessions.
\end{itemize}
I will analyse episodes by linking talk, action and system records, and
analyse interviews thematically. The framework is Lucy Suchman's account
of how agency is figured at the human--machine interface, complemented by
improvisation studies' attention to listening and response.

The project fits Technology in Practice, with Media \& Culture as a minor. I
hope to work with Dr.\ Rebecca Caines, whose research on improvisation as a
research methodology, socially responsive creative technologies, and
generative AI in research-creation fits this approach.

My background may appear distant from communication and culture, so I want
to explain how it prepares me. I hold a Bachelor of Architecture from Tabriz
Islamic Art University and a Master of Landscape Architecture from Istanbul
Technical University. There I studied how people perceive virtual
environments, work published at eCAADe 2023. I am completing a thesis-based
master's degree in Interactive Arts and Technology at Simon Fraser
University, where I built and evaluated multiplayer Unity systems, including
ambient cues that let remote teammates perceive one another's heart-rate
state without emotion labels; two manuscripts from this work are under
revision. Participants made acceptance of that display conditional on
controlling who could see it, which led me to ask who controls how a shared
system interprets people. I have conducted ethics-approved mixed-methods
studies, worked as a teaching assistant in animation and interactive media,
and served on SFU's Senate Graduate Studies Committee.

\end{document}
